% Standalone section body. Requires amsmath, booktabs, graphicx, and natbib.
% This is a prospective test-evaluation draft, not a report of completed tests.
% Each AFResult key identifies an unmeasured result; see EXPERIMENTS_DRAFT_NOTES.md.
% Override individual fields with \AFSetResult{key}{value} before this input.
\providecommand{\AFSetResult}[2]{\expandafter\def\csname AFresult:#1\endcsname{#2}}
\providecommand{\AFResult}[1]{\ifcsname AFresult:#1\endcsname\csname AFresult:#1\endcsname\else\textbf{TBD}\fi}

\section{Experiments}
\label{sec:experiments}

We evaluate predictive accuracy, satisfaction of explicit mechanism requirements,
and the contributions of scientific feedback and model simplification. We also
test sensitivity to familiar terminology and canonical dynamics. Predictive
accuracy and mechanism compliance are reported separately: a model can predict
well while representing the wrong mechanism.

\subsection{Experimental setup}
\label{sec:exp_setup}

\paragraph{Benchmarks.}
We use nine benchmark cases from three families. The Dalla Man glucose--insulin
system~\citep{dalla2007meal} supplies T1-easy, T1-hard, T2-easy, and T2-hard,
which vary the requested mechanisms and available observations. T1-easy also
includes obfuscated, perturbed, and jointly obfuscated--perturbed variants,
giving seven Dalla Man cases. We add a named controlled stirred-tank reactor
(CSTR) and a procedurally specified dynamical device. These are nine evaluation
cases, not nine independent physical systems. The appendix specifies the public
channels, outputs, and mechanism requirements of each case.

\paragraph{Protocol.}
Parameters are fitted on training trajectories; validation selects the retained
model and pruning endpoint. Models, parameters, and causal initializers are
frozen before test evaluation. Test trajectories are excluded from proposal,
critique, fitting, pruning, and selection. We use GPT-OSS-20B as proposer and
GPT-OSS-120B as advisory scientific critic. Each Autoformalism run has an initial
construction and fourteen revision visits, with two seeds per case and prompt
variant. We report both \emph{Full} and \emph{Brief-only}: the latter omits the
trajectory summary from the initial construction prompt, while subsequent
fitting and revision follow the same protocol. Neither variant is selected
using test results.

\paragraph{Metrics and aggregation.}
Our primary predictive metric is test free-rollout NMSE, normalized by training
output variances and averaged over target channels. Predicted targets are not
reset to observations after initialization; permitted auxiliary channels remain
available under the task's observation protocol. For \emph{formal mechanism
compliance}, we evaluate explicit predicates on the assembled equations after
substituting fitted parameters, independently of proposer annotations and LLM
scores. Examples include a generated input-to-target path, an input-driven
dynamic memory contributing to the target, and distinct required contributions
to an output equation. Each obligation is pass, fail, or unresolved; unresolved
obligations stay in the denominator. Passing these checks does not certify
general scientific correctness. Fitted-activity probes and response-shape
diagnostics are reported separately on their supported cases. Complexity counts
additive terms, with states, parameters, and expression size in the appendix.
Cost includes proposer and critic calls, including retries.

We report medians and unscaled median absolute deviations (MAD),
$\operatorname{MAD}(z)=\operatorname{median}_r
|z_r-\operatorname{median}_{r'}z_{r'}|$.
For each method and case we first summarize seeds; the overall result is the
median and MAD of the nine case medians, giving each case equal weight.
Full and Brief-only remain separate. Cross-case MAD describes benchmark
heterogeneity, not a confidence interval. Completion coverage accompanies these summaries;
failed and unavailable evaluations are identified rather than silently discarded.
The appendix gives per-case results and a family-balanced sensitivity analysis.

\subsection{Comparison with external baselines}
\label{sec:exp_baselines}

We compare with SINDy~\citep{brunton2016discovering},
PySR~\citep{cranmer2023interpretable}, D3~\citep{holt2024data}, and an end-to-end
GPT-5.6 Sol agent. All methods use the same permitted channels and data splits;
methods supporting language additionally receive the public task specification.
SINDy and PySR use their native numerical discovery procedures. The agent
generates and fits its own executable model; we evaluate its submitted constants
without an additional Autoformalism refit. D3 is evaluated through its native
discrete increment rollout, whereas continuous-time candidates use ODE rollouts.
Native training and validation-selection objectives, repetition counts, and
budgets are disclosed in the appendix; a common evaluation protocol does not
imply identical optimization budgets.

\begin{table}[t]
\centering
\small
\setlength{\tabcolsep}{4pt}
\caption{External comparison on the frozen test split. NMSE, terms, and tokens
are case-macro median [MAD]. P and U are case-macro median percentages of passed
and unresolved formal obligations; their MADs and per-case counts appear in the
appendix. Coverage is evaluated/planned runs, with failures itemized. D3 uses
native discrete rollouts. Every \textbf{TBD} awaits a measured result.}
\label{tab:exp_baselines}
\begin{tabular}{lccccc}
\toprule
Method & NMSE $\downarrow$ & P $\uparrow$ / U & Terms $\downarrow$ &
\shortstack{Tokens\\($10^3$)} & Coverage \\
\midrule
Ours (Full) & \AFResult{base.full.error} & \AFResult{base.full.mechanism} & \AFResult{base.full.terms} & \AFResult{base.full.tokens} & \AFResult{base.full.coverage} \\
Ours (Brief-only) & \AFResult{base.brief.error} & \AFResult{base.brief.mechanism} & \AFResult{base.brief.terms} & \AFResult{base.brief.tokens} & \AFResult{base.brief.coverage} \\
SINDy & \AFResult{base.sindy.error} & \AFResult{base.sindy.mechanism} & \AFResult{base.sindy.terms} & --- & \AFResult{base.sindy.coverage} \\
PySR & \AFResult{base.pysr.error} & \AFResult{base.pysr.mechanism} & \AFResult{base.pysr.terms} & --- & \AFResult{base.pysr.coverage} \\
D3 & \AFResult{base.d3.error} & \AFResult{base.d3.mechanism} & \AFResult{base.d3.terms} & \AFResult{base.d3.tokens} & \AFResult{base.d3.coverage} \\
GPT-5.6 Sol agent & \AFResult{base.agent.error} & \AFResult{base.agent.mechanism} & \AFResult{base.agent.terms} & \AFResult{base.agent.tokens} & \AFResult{base.agent.coverage} \\
\bottomrule
\end{tabular}
\end{table}

Table~\ref{tab:exp_baselines} compares accuracy, formal compliance, complexity,
and cost on a common benchmark roster. Across matched cases, the ratios of each
baseline's median NMSE to ours and the number of cases favoring each method are
\AFResult{base.paired_comparison_sentence}.
The corresponding compliance--accuracy trade-off is
\AFResult{base.mechanism_comparison_sentence}.
These comparisons distinguish predictive success from satisfaction of the
specified model requirements; no LLM preference score supplies the evaluation
label.

\subsection{Ablation study}
\label{sec:exp_ablations}

\paragraph{Scientific critic and deterministic verifier.}
We run a $2\times2$ ablation on all nine cases: both components, verifier only,
critic only, and neither. The critic supplies scientific revision advice but
does not select or veto models. Removing the verifier disables scientific
admission, repair-feedback, and pruning gates, while preserving executable-model
checks such as grammar, defined targets, causal initialization, and absence of
algebraic cycles. Every arm retains the same specification, fitting budget, and
revision allowance; the same independent equation assessment is applied to all
arms. Thus, this experiment separates the effect of scientific advice from
enforcement of formal requirements.

\begin{table}[t]
\centering
\small
\setlength{\tabcolsep}{5pt}
\caption{Critic/verifier ablation on all nine cases. C: scientific critic;
V: deterministic scientific verifier. Each entry gives case-macro median [MAD].
Formal compliance counts unresolved obligations in its denominator; unresolved
rates and completion coverage are reported alongside the per-case results.}
\label{tab:exp_critic_verifier}
\begin{tabular}{lccccc}
\toprule
Prompt & C & V & Test NMSE $\downarrow$ & Formal pass (\%) $\uparrow$ & Tokens ($10^3$) \\
\midrule
Full & On & On & \AFResult{cv.full.both.error} & \AFResult{cv.full.both.pass} & \AFResult{cv.full.both.tokens} \\
Full & Off & On & \AFResult{cv.full.verifier.error} & \AFResult{cv.full.verifier.pass} & \AFResult{cv.full.verifier.tokens} \\
Full & On & Off & \AFResult{cv.full.critic.error} & \AFResult{cv.full.critic.pass} & \AFResult{cv.full.critic.tokens} \\
Full & Off & Off & \AFResult{cv.full.neither.error} & \AFResult{cv.full.neither.pass} & \AFResult{cv.full.neither.tokens} \\
\midrule
Brief-only & On & On & \AFResult{cv.brief.both.error} & \AFResult{cv.brief.both.pass} & \AFResult{cv.brief.both.tokens} \\
Brief-only & Off & On & \AFResult{cv.brief.verifier.error} & \AFResult{cv.brief.verifier.pass} & \AFResult{cv.brief.verifier.tokens} \\
Brief-only & On & Off & \AFResult{cv.brief.critic.error} & \AFResult{cv.brief.critic.pass} & \AFResult{cv.brief.critic.tokens} \\
Brief-only & Off & Off & \AFResult{cv.brief.neither.error} & \AFResult{cv.brief.neither.pass} & \AFResult{cv.brief.neither.tokens} \\
\bottomrule
\end{tabular}
\end{table}

Table~\ref{tab:exp_critic_verifier} shows
\AFResult{cv.interpretation_sentence}.
We quantify each component by matched case/prompt/seed differences with the
other component fixed, rather than by comparing unmatched pools of successful
runs. Additional critic cost is reported explicitly.

\paragraph{Shared processes and pruning.}
We isolate these components on canonical named T1-easy, T1-hard, T2-easy, and
T2-hard with both critic and verifier enabled. Disabling shared processes removes
the dedicated proposal and assembly stage; ordinary algebraic reuse remains
possible. Pruning is compared with the identical final search parent and with an
unchanged-model refit using the same additional fitting budget. Deletion is
ranked on training rollouts; validation chooses among the parent, control, and
pruned candidate using the fixed tolerance described in the method.
Relative to these matched controls, the test-error and complexity changes are
\AFResult{secondary.interpretation_sentence}.
This restricted comparison avoids a Cartesian sweep of all components. Detailed
paired results, including skipped and rejected deletions, appear in the appendix.

\subsection{How specifications guide discovery}
\label{sec:exp_specification_case}
% Insert Orion's case study here. No conclusions from that experiment are
% assumed by the baseline, ablation, or memorization draft.
\AFResult{specification.case_study_text}

\subsection{Dependence on familiar models and terminology}
\label{sec:exp_memorization}

To investigate whether success depends on familiar equations or semantic cues,
we use the T1-easy $2\times2$ design: canonical versus perturbed dynamics,
crossed with named versus obfuscated presentation
(Table~\ref{tab:exp_memorization}). Named and obfuscated versions share numerical
data up to renaming and preserve the operational task requirements. Canonical
and perturbed dynamics use matched experimental schedules but different
generating dynamics. SINDy and PySR provide numerical controls because they do
not consume the textual specification.

\begin{table}[t]
\centering
\small
\setlength{\tabcolsep}{5pt}
\caption{T1-easy controls for dependence on familiar cues. Test NMSE is median
[MAD] across repetitions within each cell. Ratios are computed within matched
repetitions before aggregation. Presentation changes preserve numerical data;
changing dynamics also changes the numerical identification problem.}
\label{tab:exp_memorization}
\begin{tabular}{lcccc}
\toprule
& \multicolumn{2}{c}{Canonical dynamics} & \multicolumn{2}{c}{Perturbed dynamics} \\
\cmidrule(lr){2-3}\cmidrule(lr){4-5}
Method & Named & Obfuscated & Named & Obfuscated \\
\midrule
Ours (Full) & \AFResult{mem.full.cn} & \AFResult{mem.full.co} & \AFResult{mem.full.pn} & \AFResult{mem.full.po} \\
Ours (Brief-only) & \AFResult{mem.brief.cn} & \AFResult{mem.brief.co} & \AFResult{mem.brief.pn} & \AFResult{mem.brief.po} \\
SINDy & \AFResult{mem.sindy.cn} & \AFResult{mem.sindy.co} & \AFResult{mem.sindy.pn} & \AFResult{mem.sindy.po} \\
PySR & \AFResult{mem.pysr.cn} & \AFResult{mem.pysr.co} & \AFResult{mem.pysr.pn} & \AFResult{mem.pysr.po} \\
D3 & \AFResult{mem.d3.cn} & \AFResult{mem.d3.co} & \AFResult{mem.d3.pn} & \AFResult{mem.d3.po} \\
GPT-5.6 Sol agent & \AFResult{mem.agent.cn} & \AFResult{mem.agent.co} & \AFResult{mem.agent.pn} & \AFResult{mem.agent.po} \\
\bottomrule
\end{tabular}
\end{table}

For each method, we summarize the paired ratios
$E_{\mathrm{obfuscated}}/E_{\mathrm{named}}$ within a dynamics condition and
$E_{\mathrm{perturbed}}/E_{\mathrm{canonical}}$ within a presentation, where
$E$ is test NMSE. The resulting sensitivity pattern is
\AFResult{mem.interpretation_sentence}.
A naming gap measures dependence on semantic cues; robustness to perturbed
dynamics provides evidence against success through unchanged canonical-equation
retrieval alone. Neither result identifies training-set memorization: semantic
information, optimization difficulty, and model misspecification can also
produce these gaps. We therefore interpret this experiment as a controlled
test of reliance on familiar models, rather than proof of data contamination.
